\section{Нормальная форма Хомского}
	\tikzsetfigurename{CFG_CNF_}
\label{sec:CNF}

\begin{definition}[Нормальная форма Хомского]
    Контекстно-свободная грамматика $\langle \Sigma, N, P, S\rangle$ находится в \emph{Нормальной Форме Хомского}, если она содержит только правила следующего вида:
    \begin{itemize}
        \item $A \to B C$, где $A, B, C \in N$, а стартовый нетерминал $S$ не содержится в правой части правила.
        \item $A \to a$, где $A \in N$, $a \in \Sigma$.
        \item $S \to \varepsilon$: только из стартового нетерминала выводима пустая строка.
    \end{itemize}
\end{definition}

\begin{theorem}
    Любую КС грамматику можно преобразовать в НФХ.
\end{theorem}

\begin{proof}
    Алгоритм преобразования в НФХ состоит из следующих шагов:
    \begin{itemize}
        \item Замена неодиночных терминалов
        \item Удаление длинных правил
        \item Удаление $\varepsilon$-правил
        \item Удаление цепных правил
        \item Удаление бесполезных нетерминалов
    \end{itemize}
    То, что каждый из этих шагов преобразует грамматику к эквивалентной, при этом является алгоритмом, доказано в следующих леммах.
\end{proof}

\begin{lemma}
    Для любой КС-грамматики можно построить эквивалентную, которая не содержит правила с неодиночными терминалами.
\end{lemma}

\begin{proof}
    Каждое правило $A \to B_0 B_1 \dots B_k$, $k \geq 1$ заменить на множество правил, где $C_i$~--- новый нетерминал:
    \begin{align*}
        A   & \to C_0 C_1 \dots C_k \\
        C_0 & \to B_0               \\
        C_1 & \to B_1               \\
            & \dots                 \\
        C_k & \to B_k
    \end{align*}
\end{proof}

\begin{lemma}
    Для любой КС-грамматики можно построить эквивалентную, которая не содержит правил длины больше 2.
\end{lemma}

\begin{proof}
    Каждое правило $A \to B_0 B_1 \dots B_k$, $k \geq 2$ заменить на множество правил:
    \begin{align*}
        A       & \to B_0 C_0         \\
        C_0     & \to B_1 C_1         \\
                & \dots               \\
        C_{k-3} & \to B_{k-2} C_{k-2} \\
        C_{k-2} & \to B_{k-1} B_k
    \end{align*}
\end{proof}

\begin{lemma}
    Для любой КС-грамматики можно построить эквивалентную, не содержащую $\varepsilon$-правил.
\end{lemma}

\begin{proof}
    Рекурсивно определим $\varepsilon$-правила:
    \begin{itemize}
        \item $A \to \varepsilon$~--- $\varepsilon$-правило
        \item $A \to B_0 \dots B_k$~--- $\varepsilon$-правило, если $\forall i$: $B_i$~--- $\varepsilon$-правило.
    \end{itemize}
    Каждое правило $A \to B_0 B_1 \dots B_k$ заменяем на множество правил, где каждое $\varepsilon$-правило удалено во всех возможных комбинациях.
\end{proof}

Стоит отметить, что порядок выполнения шагов удаления длинных и $\varepsilon$-правил влияет на размер результирующей грамматики.
При удалении $\varepsilon$-правил для каждого правила $A \to B_0 \dots B_k$ порождаются все комбинации вхождений nullable-нетерминалов, что даёт до $2^m$ вариантов, где $m$~--- количество nullable-нетерминалов в правой части.
Если сперва удалить $\varepsilon$-правила, а затем~--- длинные, то на момент удаления $\varepsilon$-правил $k$ (а значит, и $m$) может быть произвольно большим, и рост грамматики экспоненциален.
Если же сперва удалить длинные правила, а затем~--- $\varepsilon$-правила, то к моменту $\varepsilon$-удаления длина любого правила не превышает двух, и каждое правило даёт не более четырёх вариантов~--- рост ограничен заранее известной константой.
Этим объясняется разнобой в оценках размера грамматики после приведения к НФХ, встречающийся в различных источниках.

\begin{lemma}
    Для любой КС-грамматики можно построить эквивалентную, не содержащую цепные правила.
\end{lemma}

\begin{proof}
    \emph{Цепное правило}~--- правило вида $A \to B\text{, где } A, B \in N$.
    \emph{Цепная пара}~--- упорядоченная пара $(A,B)$, в которой $A\derives B$, используя только цепные правила.

    Алгоритм:
    \begin{enumerate}
        \item Найти все цепные пары в грамматике $G$.
              Найти все цепные пары можно по индукции:
              Базис: $(A,A)$~--- цепная пара для любого нетерминала, так как $A\derives A$ за ноль шагов.
              Индукция: Если пара $(A,B_0)$~--- цепная, и есть правило $B_0 \to B_1$, то $(A,B_1)$~--- цепная пара.
        \item Для каждой цепной пары $(A,B)$ добавить в грамматику $G'$ все правила вида $A \to a$, где $B \to a$~--- нецепное правило из $G$.
        \item Удалить все цепные правила
    \end{enumerate}
    Пусть $G$~--- контекстно-свободная грамматика. $G'$~--- грамматика, полученная в результате применения алгоритма к $G$. Тогда $L(G)=L(G')$.
\end{proof}

\begin{definition}[Порождающий и непорождающий нетерминалы]
    Нетерминал $A$ называется \emph{порождающим}, если из него может быть выведена конечная терминальная цепочка.
    Иначе он называется \emph{непорождающим}.
\end{definition}

\begin{lemma}
    Можно удалить все бесполезные (непорождающие) нетерминалы.
\end{lemma}

\begin{proof}
    После удаления из грамматики правил, содержащих непорождающие нетерминалы, язык не изменится, так как непорождающие нетерминалы по определению не могли участвовать в выводе какого-либо слова.

    Алгоритм нахождения порождающих нетерминалов:
    \begin{enumerate}
        \item Множество порождающих нетерминалов пустое.
        \item Найти правила, не содержащие нетерминалов в правых частях и добавить нетерминалы, встречающихся в левых частях таких правил, в множество.
        \item Если найдено такое правило, что все нетерминалы, стоящие в его правой части, уже входят в множество, то добавить в множество нетерминалы, стоящие в его левой части.
        \item Повторить предыдущий шаг, если множество порождающих нетерминалов изменилось.
    \end{enumerate}
    В результате получаем множество всех порождающих нетерминалов грамматики, а все нетерминалы, не попавшие в него, являются непорождающими.
    Их можно удалить.
\end{proof}

Наряду с непорождающими, существуют также \emph{недостижимые} нетерминалы~--- те, которые не могут появиться ни в одной сентенциальной форме ни при каком выводе из стартового нетерминала $S$.

\begin{definition}[Недостижимый нетерминал]
    Нетерминал $A \in N$ называется \emph{недостижимым}, если не существует вывода $S \derives \alpha A \beta$ ни для каких $\alpha, \beta \in (\Sigma \cup N)^*$.
\end{definition}

Алгоритм поиска достижимых нетерминалов строится аналогично поиску порождающих~--- с использованием рабочего множества, но в направлении <<сверху вниз>>.
Изначально множество достижимых содержит $\{S\}$.
Далее, пока множество меняется, для каждого достижимого нетерминала $A$ и каждого правила $A \to X_1 \dots X_k$ все нетерминалы среди $X_i$ добавляются в множество достижимых.
Нетерминалы, не попавшие в итоговое множество, являются недостижимыми и могут быть удалены вместе со всеми правилами, в которых они участвуют.

Непорождающие и недостижимые нетерминалы вместе называются \emph{бесполезными}.
Важно, что порядок их удаления имеет значение.
Если сперва удалить недостижимые, а затем~--- непорождающие, одной итерации достаточно.
Если же сперва удалить непорождающие, то удаление некоторых правил может сделать отдельные нетерминалы недостижимыми, и потребуется повторная чистка.
\begin{example}
В качестве примера рассмотрим грамматику
\[
    S \to A \mid B,\qquad A \to a,\qquad B \to b,\qquad C \to C.
\]
Нетерминал $C$ непорождающий (из него ничего не выводится), и если удалить непорождающие первыми, правило $C \to C$ исчезнет, но $C$ по-прежнему останется в множестве нетерминалов~--- потребуется второй проход для удаления недостижимых.
\end{example}

\begin{example}
    Приведем в Нормальную Форму Хомского однозначную грамматику правильных скобочных последовательностей: $S \to a S b S \mid \varepsilon$

    Первым шагом добавим новый нетерминал и сделаем его стартовым:
    \begin{align*}
        S_0 & \to S                              \\
        S   & \to a \ S \ b \ S \mid \varepsilon
    \end{align*}
    Заменим все терминалы на новые нетерминалы:
    \begin{align*}
        S_0 & \to S                              & L & \to a \\
        S   & \to L \ S \ R \ S \mid \varepsilon & R & \to b
    \end{align*}
    Избавимся от длинных правил:
    \begin{align*}
        S_0 & \to S                       & S'' & \to R \ S \\
        S   & \to L \ S' \mid \varepsilon & L   & \to a     \\
        S'  & \to S \ S''                 & R   & \to b
    \end{align*}
    Избавимся от $\varepsilon$-продукций:
    \begin{align*}
        S_0 & \to S \mid \varepsilon & S'' & \to R   \mid R \ S \\
        S   & \to L \ S'             & L   & \to a              \\
        S'  & \to S'' \mid S \ S''   & R   & \to b
    \end{align*}
    Избавимся от цепных правил:
    \begin{align*}
        S_0 & \to L \ S' \mid \varepsilon   & S'' & \to b \mid R \ S \\
        S   & \to L \ S'                    & L   & \to a            \\
        S'  & \to b \mid R \ S \mid S \ S'' & R   & \to b
    \end{align*}
\end{example}

\begin{definition}[Ослабленная нормальная форма Хомского (ОНФХ)]
    \label{def:wCNF}
    Контекстно-свободная грамматика $\langle \Sigma, N, P, S\rangle$ находится в \emph{ослабленной Нормальной Форме Хомского}, если она содержит только правила следующего вида:
    \begin{itemize}
        \item $A \to B C$, где $A, B, C \in N$;
        \item $A \to a$, где $A \in N$, $a \in \Sigma$;
        \item $A \to \varepsilon$, где $A \in N$.
    \end{itemize}
\end{definition}

То есть ослабленная НФХ отличается от НФХ тем, что:
\begin{enumerate}
    \item $\varepsilon$ может выводиться из любого нетерминала;
    \item $S$ может появляться в правых частях правил.
\end{enumerate}

Алгоритм построения ОНФХ легко получается из алгоритма построения НФХ: поскольку ограничений меньше, часть шагов просто не нужна~--- не требуется добавлять новый стартовый символ (так как $S$ может появляться в правых частях), не требуется специальная обработка $\varepsilon$-правил (так как $\varepsilon$ разрешён из любого нетерминала).
Оставшиеся шаги~--- замена неодиночных терминалов, укорачивание длинных правил, удаление цепных правил~--- выполняются так же, как и для НФХ.
